\chapter{Problem Statement}
\setcounter{equation}{0}
\section{Problem Addressed}
Phishing URL detectors may achieve strong results when training and testing use the same dataset, yet perform less reliably on URLs from other sources. This limits confidence in deploying such models as browser protection. URL features can also vary in availability and predictive value across datasets, making feature selection a key part of generalisation.

\noindent\textbf{Problem statement:} How can an XGBoost-based URL phishing detector be designed and evaluated to balance detection quality, computational efficiency, and generalisation across distinct URL datasets?

\section{Research Questions}
The paper investigates three questions: (1) whether XGBoost achieves cross-dataset generalisation compared with baseline methods; (2) how feature-importance stability relates to cross-dataset performance; and (3) how the computational efficiency of XGBoost compares with neural networks for real-time deployment.\cite{bib1}

\section{Scope of the Seminar}
This report reviews the published X-PHIDE method, its datasets, feature extraction and selection, Bayesian hyperparameter optimisation, browser plug-in workflow, evaluation results, strengths, limitations, and possible future work. The numerical results are those reported by the original authors; no independent training or replication is claimed.
